\section{Interupsi DOS}
\label{sec:interupsi-dos}

DOS menyediakan layanan melalui \code{INT 21h}. Nomor fungsi
diletakkan pada \reg{AH}, sedangkan argumen lain diletakkan pada
register sesuai fungsi. Layanan DOS dipakai untuk mencetak karakter,
mencetak string, membaca masukan, dan mengakhiri program.

Bagian ini diringkas dari \cite{gunawan2019}, Bagian 18,
hlm.~179--190, dan \cite{sto2001}, BAB VII, hlm.~27--31, BAB XII,
hlm.~63--67, serta BAB XV, hlm.~78--84.

\subsection{Fungsi INT 21h yang Sering Dipakai}

Fungsi utama \code{INT 21h} dirangkum pada
Tabel~\ref{tab:int21h}.

\begin{table}[H]
    \centering
    \caption{Fungsi \code{INT 21h} yang sering dipakai.}
    \label{tab:int21h}
    \small
    \begin{tabular}{@{}p{1.6cm}p{5.2cm}p{6.6cm}@{}}
        \toprule
        \textbf{AH} & \textbf{Fungsi} & \textbf{Register dan hasil} \\
        \midrule
        \hex{01} & Membaca satu karakter dengan \eng{echo} &
        Hasil di \reg{AL} \\
        \addlinespace
        \hex{02} & Mencetak satu karakter &
        \reg{DL} = karakter yang dicetak \\
        \addlinespace
        \hex{09} & Mencetak string berakhir \code{'\$'} &
        \reg{DS}:\reg{DX} = alamat string \\
        \addlinespace
        \hex{0A} & Membaca string berpenyangga &
        \reg{DS}:\reg{DX} = alamat penyangga \\
        \addlinespace
        \hex{4C} & Mengakhiri program &
        \reg{AL} = kode keluar \\
        \bottomrule
    \end{tabular}
\end{table}

\subsection{Mencetak Karakter dan String}

Fungsi \hex{02} mencetak satu karakter yang berada di \reg{DL}.

\begin{lstlisting}[language=MASM, numbers=none]
MOV AH, 02h         ; fungsi cetak karakter
MOV DL, 'A'         ; karakter yang dicetak
INT 21h
\end{lstlisting}

Fungsi \hex{09} mencetak string yang ditunjuk \reg{DS}:\reg{DX}.
String harus diakhiri karakter \code{'\$'}. Karakter \code{'\$'}
tidak ikut tercetak.

\begin{lstlisting}[language=MASM, numbers=none]
.DATA
    pesan DB 'Halo, dunia!', 0Dh, 0Ah, '$'

.CODE
    MOV AH, 09h
    LEA DX, pesan
    INT 21h
\end{lstlisting}

\subsection{Membaca Masukan}

Fungsi \hex{01} membaca satu karakter dan menampilkannya kembali
(\eng{echo}). Hasilnya berada di \reg{AL}.

\begin{lstlisting}[language=MASM, numbers=none]
MOV AH, 01h         ; baca satu karakter
INT 21h             ; hasil di AL
\end{lstlisting}

Fungsi \hex{0A} membaca satu baris masukan ke dalam penyangga.
Penyangga memiliki susunan khusus:

\begin{itemize}
    \item byte 0: kapasitas maksimum, diisi programmer;
    \item byte 1: jumlah karakter yang benar-benar diketik, diisi
          DOS;
    \item byte 2 dan seterusnya: karakter masukan.
\end{itemize}

\begin{lstlisting}[language=MASM, numbers=none]
.DATA
    nama DB 20          ; kapasitas maksimum
         DB 0           ; jumlah terisi (diisi DOS)
         DB 20 DUP(0)   ; isi masukan

.CODE
    MOV AH, 0Ah
    LEA DX, nama
    INT 21h
\end{lstlisting}

Setelah fungsi \hex{0A} selesai, byte 1 berisi jumlah karakter.
Program dapat menambahkan \code{'\$'} setelah karakter terakhir agar
isi penyangga dapat dicetak dengan fungsi \hex{09}.

\subsection{Mengakhiri Program}

Ada dua cara umum mengakhiri program DOS.

\begin{itemize}
    \item \code{INT 20h} mengakhiri program \file{.COM} dan
          mengembalikan kendali ke DOS.
    \item \code{INT 21h} fungsi \hex{4C} mengakhiri program
          \file{.EXE} atau \file{.COM} dengan kode keluar pada
          \reg{AL}.
\end{itemize}

\begin{lstlisting}[language=MASM, numbers=none]
MOV AH, 4Ch         ; fungsi terminate
MOV AL, 00h         ; kode keluar 0
INT 21h
\end{lstlisting}

\begin{infobox}
\vspace{1.1em}
Fungsi \hex{4C} lebih umum dipakai karena dapat mengembalikan kode
keluar. Kode \hex{00} berarti program selesai tanpa kesalahan,
sedangkan kode bukan nol dapat dipakai untuk menandai kesalahan.
\end{infobox}

\begin{peringatanbox}
\vspace{1.1em}
Tiga kesalahan yang sering terjadi pada layanan DOS:
\begin{enumerate}
    \item string untuk fungsi \hex{09} tidak diakhiri \code{'\$'};
    \item penyangga fungsi \hex{0A} tidak menyediakan byte kapasitas
          dan byte jumlah;
    \item lupa mengakhiri program sehingga kendali tidak kembali ke
          DOS.
\end{enumerate}
\end{peringatanbox}
